%=====================================================================
\section{Assumptions}
\label{sec:assumptions}

The assumptions fall into four groups: those inherited from formulation 3, those of the simulator (environment), those of learning, and those of the test-time search. Each is stated with the reason it was made and what happens if it fails. Where an assumption is in fact a provable property, the proof is in Section~\ref{sec:mdp}.

\subsection{Inherited from formulation 3}

The RL method solves exactly formulation 3; it adds no simplification of the model. In particular it inherits:
\begin{enumerate}[label=F\arabic*.]
  \item All requests of the period are known at time 0 (static problem); request $i$ cannot start before $r_i$.
  \item Each EV is served by exactly one gun or one van, one EV at a time per resource, without pre-emption; charging power is constant, so $\tau^\Fcs_i$ and $\tau^\Mcs_i$ are known.
  \item The income index $\theta_i$ and therefore $\mu_i$, $c^{a}_i$, $W_i$ are known numbers per EV (the code uses straight lines in $\theta$, \code{code/f3\_params.json}; the model needs only their monotonicity).
  \item Vans are identical (same $B$, $Q$, $c^\Fcs$), start and end at one depot, and are not recharged during the period.
  \item Distances are road km $=$ detour factor $\times$ km per grid unit $\times$ straight-line distance, and travel times are proportional to distances (one speed). Hence distances and times satisfy the \textbf{triangle inequality}, which several proofs below use.
\end{enumerate}

\subsection{Environment (simulator) assumptions}

\begin{assumption}[Earliest-start timing]\label{as:timing}
Once it is fixed which resource serves each EV and in which order, every session starts as early as possible: $S_i=\max(r_i+t_{if(g)},\,a_g)$ on a gun, $S_i=\max(r_i,\,a_m+t_{\ell_m i})$ on a van, where $a$ is the time the resource becomes free and $\ell_m$ the van's last node.

\textbf{Why.} With assignment and sequences fixed, the remaining problem in the start times is a linear program whose constraints are lower bounds of the form $S_j\ge S_i+\text{const}$ or $S_j\ge\text{const}$, plus upper bounds ($q^\Fcs_i\le W_i$, $C_i\le T$, return by $T$), and whose cost is non-decreasing in every $S_i$. The componentwise smallest schedule is therefore optimal whenever any feasible schedule exists (Proposition~\ref{prop:timing}).

\textbf{If it fails.} It cannot fail under F1--F5. It would fail if waiting were rewarded (e.g.\ a cheaper tariff later in the period), which formulation 3 does not have.
\end{assumption}

\begin{assumption}[Append-only construction]\label{as:append}
An action appends an EV to the \emph{end} of a resource's current sequence; the construction never inserts in the middle and never undoes a decision.

\textbf{Why.} It makes each step cheap and keeps the timing of already placed EVs fixed, so the reward of a step is final. It loses nothing: every feasible plan of formulation 3 can be produced by some order of appends (Proposition~\ref{prop:complete}).

\textbf{If it fails.} Not a restriction of the plan space, only of the order in which the policy must find it: a plan whose sequences the policy ``begins wrongly'' cannot be repaired within the episode. Sampling many episodes and active search compensate.
\end{assumption}

\begin{assumption}[A dead EV stays dead]\label{as:dead}
If an unplaced EV has no feasible (EV, resource) pair at some step, it has none at any later step; it is dropped at once, counted, and the episode continues with the others.

\textbf{Why.} Gun free times only grow, van batteries only shrink, van seats only fill, and by the triangle inequality a van that moves on can only get further in time and energy from any EV (Proposition~\ref{prop:dead}). Dropping at once saves useless steps and gives a clean count of failures.

\textbf{If it fails.} It cannot fail under F1--F5.
\end{assumption}

\begin{assumption}[Lossless symmetry reduction]\label{as:sym}
Of several guns of one station that become free at the same minute, only the first is offered; of the unused vans, only the first.

\textbf{Why.} Such resources are interchangeable: exchanging their labels maps any plan onto an equally costly plan. Offering only one representative removes duplicate actions (which would split the probability of one decision over identical copies) without removing any plan up to relabelling. It is the RL counterpart of the MILP's symmetry constraints (gun order and van order).

\textbf{If it fails.} It would fail if vans or guns of one station differed (different battery, power or cost); the dataset has none.
\end{assumption}

\begin{assumption}[Exact incremental reward]\label{as:reward}
The cost increase of an action is computed with the same terms as \code{check\_plan\_f3}, including the van's way back to the depot, so the rewards of a complete episode add up exactly to $-Z$ (Proposition~\ref{prop:reward}). All times and costs are computed in 64-bit floating point.

\textbf{Why.} The policy must optimise the true objective, not a proxy. 64-bit arithmetic makes the simulator's feasibility decisions agree with the checker's (tolerance $10^{-4}$); in tests the simulator's $Z$ and the checker's $Z$ agree to $10^{-6}$ relative.
\end{assumption}

\begin{assumption}[Penalty for dropped EVs]\label{as:penalty}
The training and selection cost of a rollout is $\tilde c=Z+\lambda_{\text{drop}}\,\sigma\,|D|$, where $D$ is the set of dropped EVs, $\lambda_{\text{drop}}=10$ and $\sigma$ is the instance's cost scale: the mean, over EVs, of the cheapest feasible first option.

\textbf{Why.} A dropped EV breaks the hard rule ``every request served'', so it must cost much more than serving it. Measuring the penalty in units of a typical EV's cost ($\sigma$) makes it comparable across instances whose $Z$ differ by orders of magnitude (C2 vs R1, $n=10$ vs $n=100$).

\textbf{If it fails.} Too small: the policy might drop an awkward EV to save cost. Too large: the gradient is dominated by drops and learns little about cost. At test time selection is lexicographic (fewest drops first), so the value of $\lambda_{\text{drop}}$ does not affect which plan is reported.
\end{assumption}

\begin{assumption}[Completing an invalid plan for reporting]\label{as:complete}
If the best plan still has dropped EVs, each dropped EV is appended to the resource that finishes it earliest, ignoring the limit it violates; the plan is then scored by the checker, which lists the violations, and reported with status \code{invalid\_plan}.

\textbf{Why.} The common results format needs a resource for every EV. The plan is never presented as valid: the violations are counted and the instance is reported as ``no valid plan''.
\end{assumption}

\subsection{Learning assumptions}

\begin{assumption}[Training distribution]\label{as:dist}
The policy is trained only on random instances made with the EV-Solomon rules (Section~\ref{sec:generator}) and is tested on the 56 Solomon files it has never seen. It is assumed that the test files are typical draws of these rules.

\textbf{Why.} Within a Solomon class all files share positions, demands, ranges and $\theta$; only request times differ. Training on them would mean learning the test set. Random instances from the same rules give unlimited fresh data.

\textbf{If it fails.} If the generator misses a feature of the test files (e.g.\ the exact shape of request times in one class), the policy may be worse on that class. The RL1 results by class (Table~\ref{tab:rl1class}) are the check: the mean gap to the MILP stays between 0.6\,\% and 1.6\,\% in all six classes, so no class is badly out of distribution.
\end{assumption}

\begin{assumption}[Size and fleet coverage]\label{as:size}
Training uses the templates of the 100-EV files only (12 stations, 24 guns, 5 vans) with $n\in[10,100]$; one network is used for every $n$.

\textbf{Why.} The 560 test instances are all of that size. The attention architecture is size-free, so the same weights apply to every $n$; RL2 shifts training towards large $n$, where RL1 was weakest.

\textbf{If it fails.} For the 200-EV files (24 stations, 12 vans) the network runs, but its quality there has not been measured.
\end{assumption}

\begin{assumption}[Greedy cost rule as teacher and yardstick]\label{as:teacher}
The warm-up imitates the greedy cost rule of \code{solve\_milp\_f3} (stranded EVs first, then by request time, then higher income first; each takes its cheapest feasible resource), and validation measures the policy's cost relative to that rule.

\textbf{Why.} The rule is the MILP's own warm start: cheap, deterministic, and a reasonable plan in most instances. Imitating it gives REINFORCE a sensible start. Comparing with it on fixed instances is a fast, stable measure of progress; the MILP is far too slow to be a validation yardstick during training.

\textbf{If it fails.} A teacher bias could remain. REINFORCE after the warm-up must move away from the rule to improve; RL1's validation ratio fell from 1.01 after the warm-up to 0.969, so it did.
\end{assumption}

\begin{assumption}[Hyperparameters]\label{as:hyper}
Network sizes follow Kool et al.\ (embedding 128, 8 heads, 3 layers, feed-forward 512, clipping 10). Learning rates, batch sizes, penalty and search settings were set by hand from laptop tests; none was tuned systematically.

\textbf{If it fails.} Results may be below what tuned settings would give; Section~\ref{sec:limits} lists this as a limitation.
\end{assumption}

\subsection{Test-time assumptions}

\begin{assumption}[Map symmetries]\label{as:aug}
Coordinates are taken relative to the depot and scaled; the 8 dihedral transformations (swap $x$/$y$, flip signs) leave every distance unchanged, so all 8 inputs describe the same instance and every plan is valid for all of them.

\textbf{Why.} The network is not exactly invariant to rotation, so the 8 inputs give 8 different (but equally valid) policies: a free source of diversity for sampling and for active search.
\end{assumption}

\begin{assumption}[Selection of the reported plan]\label{as:select}
Among all plans built for an instance, the reported one is the one with the fewest dropped EVs, and among those the lowest $Z$ (key $=10^{12}|D|+Z$).

\textbf{Why.} Validity first, then cost, as in the MILP comparison.
\end{assumption}

\begin{assumption}[Active search budget]\label{as:eas}
Adapting only the instance embeddings for at most 200 iterations or 10 seconds is enough to gain most of what per-instance adaptation can give; $\lambda=0.01$ and learning rate $0.004$ fit all instances.

\textbf{Why.} The EAS paper reports that embedding-only adaptation is close to full active search at a fraction of its cost, and the laptop test on RL1's model gained about 2 percentage points of gap within 30 iterations (Section~\ref{sec:evidence}). The 10\,s cap bounds the total time of 560 instances by about 1.6 hours.
\end{assumption}

\begin{assumption}[Reproducibility]\label{as:seed}
Every instance gets its own seed, computed from the run seed and a checksum of ``file\_n''. Results therefore do not depend on the order in which instances are solved, and a resumed run gives the same plans as an uninterrupted one.
\end{assumption}
